\section{提案手法}
\label{sec:method}

\subsection{基本コンセプトと仮説}

本研究は，数値ラベル $\bm{c}$（目標値と現在の出力値の差分）を指定することで，
ベース動作データ $\bm{\xi}_\text{base}$ から所望の変位だけずれた動作をロボットに自律実行させることを目的とする．
ここで「動作データ」とは，位置・速度・トルクを含む関節状態の時系列 $\bm{\xi} \in \mathbb{R}^{T \times D}$ を指す．
ただし，$T$ は動作の全タイムステップ数，$D$ は関節の物理状態の総次元数である．

提案フレームワークでは，ILC によってラベルと実際の出力値を反復的に収束させることを想定している．
このとき，序論で述べた縮小写像の原理に基づき，ラベル $\bm{c}$ の増減に対して実際の制御出力 $\bm{v}(\bm{\xi})$ が極性を反転させることなく追従する「広域的な単調増加性」，全域で感度が一定である「傾きの均一性」，および固定の学習ゲインによる収束の妨げとならない「低い試行間推論分散」の3つの性質が同時に満たされることが，ILC 収束の前提条件となる．

本研究の仮説は，動作データ全体 $\bm{\xi}$ を直接予測するより，ベース動作データからの残差$\Delta\bm{\xi}$ だけを予測する方が，これら3つの性質が強まる，というものである．
ベース動作データ $\bm{\xi}_\text{base}$ という固定のアンカーが存在することで，NNが学習すべき変化量は小さくなり，ラベルと出力の対応関係がシンプルになると考えられる．
この仮説が成立することを，本稿の予備実験で検証する．

提案フレームワークは次の三段構成をとる．
\begin{enumerate}
  \item \textbf{残差データの収集}：Motion ReTouch \cite{Inami2025} を用いて，ベース動作データに対する修正後動作データ $\bm{\xi}_\text{corr}$ を収録し， ラベル $\bm{c}$ に対応する教師残差 $\Delta\bm{\xi}^{(\bm{c})} = \bm{\xi}^{(\bm{c})}_\text{corr} - \bm{\xi}_\text{base}$ を得る．
  \item \textbf{残差NNの学習}：ラベル $\bm{c}$ を条件として，観測から $\Delta\bm{\xi}^{(\bm{c})}$ を予測する残差NNを学習する．
  \item \textbf{ILCによる反復収束}：学習済みNNを用いてロボットを実行し，出力誤差をラベルに反映させる更新則を反復することで，目標値への収束を実現する．
\end{enumerate}

\subsection{問題設定}

記号を表\,\ref{tab:notation} に定義する．
ラベルは目標値と現在の出力値の差分として定義する：
\begin{equation}
  \bm{c} = \bm{v}^* - \bm{v}(\bm{\xi}_\text{base}),
  \label{eq:label}
\end{equation}
ここで $\bm{v}(\bm{\xi})$ は動作データ $\bm{\xi}$ を実行したときのプレース位置・最終トルクなどの評価量，$\bm{v}^*$ は目標値，$m$ はその次元数である．

\begin{table}[t]
  \centering
  \caption{記号定義}
  \label{tab:notation}
  \small
  \begin{tabular}{ll}
    \hline
    記号 & 定義 \\
    \hline
    $\bm{\xi}_\text{base} \in \mathbb{R}^{T \times D}$ & ベース動作データ \\
    $\bm{\xi}^{(\bm{c})}_\text{corr} \in \mathbb{R}^{T \times D}$ & ラベル $\bm{c}$ に対応する修正後動作データ \\
    $\Delta\bm{\xi}^{(\bm{c})} \in \mathbb{R}^{T \times D}$ & 教師残差 $:= \bm{\xi}^{(\bm{c})}_\text{corr} - \bm{\xi}_\text{base}$ \\
    $\bm{c} \in \mathbb{R}^m$ & 試行レベルのラベル \\
    $\bm{v}(\bm{\xi}) \in \mathbb{R}^m$ & 動作データ $\bm{\xi}$ の評価量 \\
    $\bm{v}^* \in \mathbb{R}^m$ & 目標値 \\
    \hline
  \end{tabular}
\end{table}

\subsection{モデルアーキテクチャ}

本研究では，Transformerデコーダに基づく模倣学習モデルを用いる．
入力は過去 $L$ ステップのフォロワアームの関節状態系列
（各ステップで角度・角速度・トルクを結合した $d_a$ 次元ベクトル），現在時刻のカメラ画像，および数値ラベル $\bm{c}$ である．
% 出力は将来 $L$ ステップの関節状態系列 $\hat{\bm{y}} \in \mathbb{R}^{L \times d_a}$ である．
出力は将来 $L$ ステップの関節状態の予測残差系列 $\Delta\hat{\bm{y}} \in \mathbb{R}^{L \times d_a}$ である．
この $\Delta\hat{\bm{y}}$ は，教師残差 $\Delta\bm{\xi}^{(\bm{c})} \in \mathbb{R}^{T \times D}$ のうちフォロワアーム成分（$d_a$ 次元）を $L$ ステップずつ予測したものに対応する．

状態系列は線形層と正弦波位置エンコーディングによりトークン化し，Transformerデコーダのクエリとして用いる．
画像は固定重みの ResNet-50 \cite{he2016resnet}で特徴マップを抽出し，Transformerエンコーダを経てメモリトークン列に変換する．
画像を使用しない場合は，学習可能なメモリトークン列をそのままキー・バリューとして用いる．
数値ラベル $\bm{c}$ は状態系列と結合して入力する．
学習は予測と教師信号のMSE損失で行う．

推論時にはAction Chunking with Transformers \cite{Zhao2023ACT} に倣ったTemporal Ensembleを適用し，
ローリングバッファに蓄積した予測の平均によって推論ノイズを低減する．
最終的な出力軌道 $\hat{\bm{y}}$ は，ベース動作データ $\bm{\xi}_\text{base}$ の対応するステップの値に対して，予測残差 $\Delta\hat{\bm{y}}$ に Temporal Ensemble をかけたものを加算することで得られる．